\section{Limitations}
\label{sec:limitations}

\textbf{Development status.} The final architecture was selected after inspecting all three
OpenCity folds during development; TrafficSH in particular was used in an earlier smoke study
(\texttt{rugd\_smoke\_v1}) that directly informed the normalized-diffusion design
(\texttt{docs/OPENCITY\_PROTOCOL.md}). All three folds are reported, but TrafficCD and TrafficHZ
are the primary confirmatory folds; TrafficSH is development evidence reported as a third fold,
not an untouched held-out geography. Target-city \emph{labels} never enter tuning or selection at
any stage (enforced in code and tested directly by
\texttt{tests/test\_opencity.py::test\_source\_statistics\_do\_not\_consult\_target}, which mutates
target-city data after fitting source statistics and asserts the fit is unchanged), but the
architecture itself was chosen with all three folds' unlabeled structure visible during
development. A genuinely unseen fourth city or an external dataset is the only way to settle
whether this design choice generalizes beyond three-city internal evidence.

\textbf{Tuning-budget asymmetry for two baselines.} As disclosed in \Cref{sec:setup}, GRU and
Graph-GRU were not tuned within this study's logged Optuna pipeline; their hyperparameters were
carried in from an earlier, unretained tuning pass. The four models that were tuned within this
study (RUGD, UIGNN-Inductive, STAEformer-Inductive, OpenCity-Scratch) remain fully evidenced and
mutually comparable; GRU/Graph-GRU's standing relative to those four should be read with this
asymmetry in mind.

\textbf{Not universally dominant.} STAEformer-Inductive wins under the two heaviest sensor-dropout
conditions (\Cref{tab:robustness}); \methodname{} is not robust to every registered condition, and
we do not claim it is. Attention over the full temporal window plausibly lets STAEformer recover
missing recent readings more gracefully than \methodname{}'s local temporal encoder plus
degree-normalized diffusion, but this mechanism has not been isolated experimentally -- a wider-
temporal-receptive-field variant of \methodname{} would be the natural follow-up ablation.

\textbf{No city-level statistical significance.} With three held-out cities, the exact sign-flip
test's minimum achievable two-sided $p$-value is 0.25 (\Cref{sec:setup}); no claim of geographic
generalization beyond these three cities is made anywhere in this paper. The paired, seed-level
Wilcoxon results (\Cref{app:paired-inference}) are strong (rank-biserial $=1.0$ against GRU,
Graph-GRU, UIGNN, and OpenCity-Scratch at both missingness levels) but characterize
seed/optimization variance on these three specific cities, not independent geographic replication.

\textbf{Uncertainty-quality asymmetry, diagnosed but not fixed.} See \Cref{sec:analysis}.

\textbf{No official baseline reproductions.} STAEformer, UIGNN, and OpenCity are clean-room task
adaptations chosen for protocol fairness within this study, not bit-identical reproductions of
their original papers; their published numbers on their own benchmarks (different labels, horizon,
and splits) are not directly comparable to anything reported here.

\textbf{Scope.} The study covers three cities, one country, and source-relative speed thresholds
rather than independently verified incident labels. Three seeds characterize optimization
variance, not geographic diversity. A publication-grade confirmatory claim needs additional cities
or an external dataset, and ideally independent or official baseline reproductions alongside these
adaptations.
